\documentclass[11pt]{article}

\usepackage[margin=1in]{geometry}
\usepackage{amsmath,amssymb,amsthm}
\usepackage{graphicx}
\usepackage[hidelinks]{hyperref}

\newtheorem*{theorem}{Theorem}

\newcommand{\T}{\mathcal T}

\begin{document}

\begin{center}
{\Large Literature status: the weakly open Daugavet-thickness question}

\medskip
{\small Run \texttt{fa\_banach\_001}; source arXiv:2005.02045}
\end{center}

\section*{Packet Status}

This packet records a literature-already-answered result, not an original
construction from the run.  The supporting paper explicitly identifies the
same question from Haller--Langemets--Lima--Nadel--Rueda Zoca and gives a
negative answer.  The packet addresses only Question (b) in the
Ivakhno-related list from the source paper.

\section*{Source Question}

The source is Rainis Haller, Johann Langemets, Vegard Lima, Rihhard Nadel, and
Abraham Rueda Zoca, \emph{On Daugavet indices of thickness},
arXiv:2005.02045.  In the final section, the authors ask:

\begin{center}
\includegraphics[width=\textwidth]{figures/open_problem_crop.png}
\end{center}

The part relevant here is Question (b):
\[
  \T(X)\ge 1
  \quad\Longrightarrow?\quad
  \text{every nonempty relatively weakly open subset of }B_X
  \text{ has diameter }2.
\]
Remark 3.9 of the same source paper says that Question (b) remains open after
their \(JT_\infty\)-predual example.

\section*{Supporting Literature}

The supporting paper is Abraham Rueda Zoca, \emph{Diameter, radius and
Daugavet index of thickness of slices in Banach spaces}, arXiv:2306.01467.
Its introduction states that the question whether \(\T(X)\ge 1\) implies that
every nonempty relatively weakly open subset of \(B_X\) has diameter two
remained open in \cite[Remark 3.9]{HLLNR}, and that this is precisely one of
the questions addressed in the note.  Its final Remark 4.4 records the answer:

\begin{center}
\includegraphics[width=\textwidth]{figures/supporting_answer_crop.png}
\end{center}

\section*{Idea of the Proof}

The source question asks whether the threshold \(\T(X)\ge 1\) forces all
relatively weakly open subsets of the unit ball to be diameter two.  Rueda
Zoca's later construction starts from a Banach space with the Daugavet
property and changes the norm.  The new norm keeps the convex-combination
thickness index very large, so in particular \(\T(X)>1\), but it also creates
a genuine slice of the new unit ball with diameter strictly smaller than two.
Since every slice is relatively weakly open, this directly violates the
conclusion asked in Question (b).

\section*{Formal Statement}

\begin{theorem}
Question (b) from arXiv:2005.02045 has a negative answer.  More strongly, for
every \(\varepsilon>0\) there is a Banach space \(X\) with
\(\T(X)>2-\varepsilon\) whose unit ball contains a nonempty relatively weakly
open subset of diameter strictly smaller than \(2\).
\end{theorem}

\begin{proof}
Rueda Zoca's Theorem 1.2 in arXiv:2306.01467 states that if \(Y\) has the
Daugavet property, then for every \(\varepsilon>0\) there is an equivalent
renorming \(|\cdot|_\varepsilon\) such that the renormed space fails the
\(r\)-big slice property, so its unit ball contains slices of diameter
strictly smaller than \(2\), and
\[
  \T^{cc}\bigl((Y,|\cdot|_\varepsilon)\bigr)\ge 2-\varepsilon
\]
after choosing the auxiliary parameters in the construction sufficiently
small.  Remark 4.4 of the supporting paper restates the consequence in the
exact language of the source question: for every \(\varepsilon>0\) there is a
Banach space \(X\) with \(\T(X)>2-\varepsilon\) and with slices of diameter
strictly smaller than \(2\).

For completeness, here is the direct implication.  The Daugavet thickness
indices satisfy
\[
  0\le \T^{cc}(X)\le \T(X)\le \T^s(X)\le 2
\]
as recalled in the supporting paper.  Thus a lower bound near \(2\) for
\(\T^{cc}(X)\) gives the same lower bound for \(\T(X)\).  Taking
\(\varepsilon<1\), the renormed space has \(\T(X)>1\).  But one of its slices
has diameter strictly smaller than \(2\).  A slice of the unit ball is a
nonempty relatively weakly open subset of the unit ball.  Therefore
\(\T(X)\ge 1\) does not force all nonempty relatively weakly open subsets of
\(B_X\) to have diameter \(2\).
\end{proof}

\section*{Limitations}

This packet covers only Question (b) from the three-part Ivakhno-related list
in the source paper.  Questions (a) and (c) were answered negatively inside
arXiv:2005.02045 itself and are therefore same-paper extraction false
positives for this run.  The separate \(\ell_1\)-sum question on
\(\ell_1\)-sum equalities for \(\T\) and \(\T^{cc}\) is not addressed here.

\section*{Literature And Duplicate Check}

The run indexes were searched for arXiv:2005.02045, Daugavet indices of
thickness, Daugavet index, weakly open subsets, Ivakhno, Rueda, and
\(\ell_1\)-sum phrases.  No existing packet or attempt covered this exact
Question (b).  The supporting paper arXiv:2306.01467 explicitly cites
\cite[Remark 3.9]{HLLNR} and states that Theorem 1.2 gives the negative
answer.  Nearby Daugavet/\(\Delta\)-point papers arXiv:2307.10647 and
arXiv:2301.04433 are adjacent but do not change the scope of this exact
literature match.

No computational checks were needed; the packet is a source-to-source
literature identification.

\section*{Human Review Recommendation}

Review as an explicit literature-already-answered packet.  The main checks
are that the later paper's \([7,\) Remark 3.9\(]\) is the Haller--Langemets--
Lima--Nadel--Rueda Zoca question shown above, and that the small-diameter
slice in the supporting theorem is indeed a relatively weakly open subset of
the renormed unit ball.

\begin{thebibliography}{9}

\bibitem{HLLNR}
R. Haller, J. Langemets, V. Lima, R. Nadel, and A. Rueda Zoca,
\emph{On Daugavet indices of thickness},
J. Funct. Anal. \textbf{280} (2021), article 108846; arXiv:2005.02045.

\bibitem{RZ2023}
A. Rueda Zoca,
\emph{Diameter, radius and Daugavet index of thickness of slices in Banach spaces},
arXiv:2306.01467.

\end{thebibliography}

\end{document}
